\SourceLeaf{10}
\begin{center}
{\Large NACHLASS.\par}
\bigskip
[ZWEI NOTIZEN ÜBER DIE AUFLÖSUNG DER CONGRUENZ
\[
xx+yy+zz\equiv 0\quad(\mathrm{mod.}\ p).]
\]
\end{center}
\begin{center}[1.]\end{center}
\noindent\hspace*{\parindent}Jede Zerlegung einer durch die Primzahl $p$ theilbaren Zahl in vier Quadrate $=aa+bb+cc+dd$ entspricht einer Auflösung der Congruenz
\[
xx+yy+zz\equiv 0\quad(\mathrm{mod.}\ p).
\]
\noindent Zusammen giebt es $pp$ solcher Auflösungen. Die Eine $x\equiv0$, $y\equiv0$, $z\equiv0$ davon weggenommen, zerfallen die übrigen $pp-1$ in $p+1$ Classen, wenn man $x=\xi$, $y=\eta$, $z=\zeta$ mit $x\equiv k\xi$, $y\equiv k\eta$, $z\equiv k\zeta$ zu Einer Classe zählt.

Es sind nämlich proportional
\[
\begin{array}{c|c|c|c}
x & aa+bb & -ad-bc & bd-ac\\
y & ac+bd & ab-cd & aa+dd\\
z & ad-bc & aa+cc & -ab-cd
\end{array}
\]
\begin{center}[2.]\end{center}
\noindent\hspace*{\parindent}Alle Auflösungen der Congruenz
\[
1+xx+yy\equiv0\quad(\mathrm{mod.}\ p)
\]
\noindent zu finden.
\clearpage
